% year: תשע"ט
% semester: א
% moed: ב
% number: 2
% points: 20
% categories: taylor
% source: exams/מבחנים/תשעט/מועד ב תשעט.pdf

\begin{question}
תהי $f(x)=\ln x$. רשמו פולינום טיילור ממעלה 3 של פונקציה $f$ בסביבת נקודה $x_0=1$. השתמשו בפולינום זה כדי לחשב קרוב ל-$\ln(1.1)$. העריכו את שגיאת הקרוב הזה בעזרת נוסחת טיילור.
\end{question}

\begin{hint}
השארית בצורת לגרנז' היא $R_3(x)=\frac{f^{(4)}(c)}{4!}(x-1)^4$ עם $c$ בין $1$ ל-$x$.
\end{hint}

\begin{solution}
\textbf{הפולינום.} $f(x)=\ln x$, $f'(x)=\frac1x$, $f''(x)=-\frac1{x^2}$, $f'''(x)=\frac2{x^3}$, $f^{(4)}(x)=-\frac6{x^4}$. בנקודה $1$: $f(1)=0$, $f'(1)=1$, $f''(1)=-1$, $f'''(1)=2$. לכן
\[ P_3(x)=(x-1)-\frac{(x-1)^2}{2}+\frac{(x-1)^3}{3}. \]
\textbf{הקירוב.} עבור $x=1.1$, $x-1=0.1$:
\[ \ln(1.1)\approx P_3(1.1)=0.1-\frac{0.01}{2}+\frac{0.001}{3}=0.1-0.005+0.000\overline{3}=0.095\overline{3}. \]
\textbf{הערכת השגיאה.} לפי נוסחת טיילור עם שארית לגרנז', קיים $c\in(1,1.1)$ כך ש-
\[ R_3(1.1)=\frac{f^{(4)}(c)}{4!}(0.1)^4=-\frac{6}{24c^4}\cdot10^{-4}=-\frac{10^{-4}}{4c^4}. \]
מכיוון ש-$c>1$, $c^4>1$ ולכן
\[ |R_3(1.1)|<\frac{10^{-4}}{4}=2.5\cdot10^{-5}. \]
כמו כן $R_3<0$, כלומר הקירוב $0.09533$ גדול מעט מהערך האמיתי. (לשם השוואה $\ln1.1=0.0953102\ldots$, והשגיאה בפועל כ-$2.3\cdot10^{-5}$.)
\end{solution}
